\section*{Chapter \ref{ch:ml:sequence-models-on-order-books} --- Sequence Models on Order Books}

\begin{solution}{exo:ml:sequence-models-on-order-books:1}
$1 + (3 - 1)(1 + 2 + 4 + 8) = 31$ snapshots; with width five, $1 + 4\times15 = 61$.
\end{solution}

\begin{solution}{exo:ml:sequence-models-on-order-books:2}
Always ``flat'': 72.7\%. Guessing in the class proportions: $0.727^2 + 0.147^2 + 0.125^2 = 56.6\%$.
\end{solution}

\begin{solution}{exo:ml:sequence-models-on-order-books:3}
$T\times T\times h = 20\times20\times32 = 12\,800$ multiply-adds; quadratic in the sequence length (and linear in the
width).
\end{solution}

\begin{solution}{exo:ml:sequence-models-on-order-books:4}
Averaging the convolution's outputs over the window weights the snapshot of twenty seconds ago as much as the current one,
although only the recent state predicts the next seconds. Read the last position instead (as the TCN and the \omterm{def:ml:sequence-models-on-order-books:attention}{transformer}
do), or flatten the time axis into the classifier.
\end{solution}

\begin{solution}{exo:ml:sequence-models-on-order-books:5}
The mean of the next ten seconds' mids moves less and more smoothly than the next move, and its first seconds continue the
current flow, which the inputs see; much of its predictability comes from that overlap. A trade earns the price change to
its exit, not an average of future mids, so a model scored on the smoothed label overstates the tradable edge; score it
on the P\&L of the trade it would drive (chapter~29).
\end{solution}

\begin{solution}{exo:ml:sequence-models-on-order-books:6}
Random splits put windows from the same day, overlapping in time, in both sets: the network is tested on near-copies of
its training windows (chapter~3). And 88\% must be compared with always predicting the majority class. Split by day,
purge the edges, and report an IC or a P\&L.
\end{solution}

\begin{solution}{exo:ml:sequence-models-on-order-books:7}
\texttt{scaling(16, name='LSTM')}: IC 0.381, against the TCN's 0.441 and the \omterm{def:ml:linear-and-regularised-baselines:logistic}{logistic regression}'s 0.568 on the same
sixteen sessions.
\end{solution}

\begin{solution}{exo:ml:sequence-models-on-order-books:8}
A causal layer with dilation $r$ at position $t$ reads $t, t - r,\dots, t - (k-1)r$: it extends the reach back by
$(k-1)r$ and never forward. Stacking layers adds the reaches: $(k-1)(1 + 2 + \dots + 2^{L-1}) = (k-1)(2^L - 1)$, plus the
position itself. By induction on the layers, every output at $t$ is a function of inputs at $t$ or earlier; the residual
connections add the same position only.
\end{solution}

\begin{solution}{pb:ml:sequence-models-on-order-books:1}
\begin{enumerate}
\item A snapshot: for each of ten levels per side, the distance from the mid in ticks and the log size in lots (40 numbers);
a window: the last 20 snapshots, ten seconds.
\item Mean mid over the next ten seconds minus the current mid, flat within a quarter tick. Test shares: 72.7\% flat,
14.7\% down, 12.5\% up.
\item Predicting ``flat'' always is right 72.7\% of the time; the chapter ranks models by the IC of $\P(\text{up}) -
\P(\text{down})$ with the forward change.
\item Windows of one session overlap and share its state: across a split they leak.
\item CNN nine snapshots per filter stack, then an average over all 20; LSTM and \omterm{def:ml:sequence-models-on-order-books:attention}{transformer} all 20; TCN 31, more than the
window.
\item CNN 0.727 and 0.060; LSTM 0.729 and 0.292; TCN 0.766 and 0.366; attention 0.775 and 0.426.
\item Accuracy 0.781, IC 0.543.
\item Standard deviations of 0.15 to 0.20 of IC for the sequence models, 0.09 for the \omterm{def:ml:linear-and-regularised-baselines:logistic}{logistic regression}.
\item After one to four \omterm{def:ml:neural-networks-for-noisy-tabular-data:backprop}{epochs}: 8\,808 windows are few for models reading 800 numbers each.
\item 0.13, 0.30 and 0.44.
\item Hardly (0.54 to 0.57): its ten features already summarise what the book says about the next seconds; it has few
parameters and little variance to reduce.
\item Keep rising past 0.57 as sessions are added, which the doubling from eight to sixteen does not yet show.
\item That simpler networks on order-flow inputs outperform most models trained directly on order books.
\item That all fifteen published models lose a large part of their performance on new data.
\item Billions of quotes pooled across many stocks.
\item \emph{Named result}: the best sequence model (attention) reaches IC 0.426 and accuracy 0.775 against the \omterm{def:ml:linear-and-regularised-baselines:logistic}{logistic
regression}'s 0.543 and 0.781, with standard deviations across sessions and seeds of 0.15 against 0.09: a gain of $-0.12$
of IC, the gap narrowing as sessions are added.
\item The \omterm{def:ml:linear-and-regularised-baselines:logistic}{logistic regression} (or boosted trees) on order-flow features; next, the TCN with the same features as extra
channels, trained on many more sessions and instruments.
\item Through chapter~2: a meta-labelled threshold on $\P(\text{up}) - \P(\text{down})$ that pays the spread only when the
expected move exceeds it, evaluated on trade P\&L.
\item More sessions and instruments to pool, and market-by-order inputs; the same split rules.
\item Much more data than the feature needs, or the feature itself as an input.
\end{enumerate}
\end{solution}

\begin{solution}{iq:ml:sequence-models-on-order-books:1}
A convolution whose output at $t$ uses only inputs at $t$ and before. A trading model must not see the future of its own
window; a centred convolution trained offline leaks the next steps into each prediction.
\iqlookfor{the definition and the look-ahead it prevents.}
\end{solution}

\begin{solution}{iq:ml:sequence-models-on-order-books:2}
$\mathsf Q = ZA_{\mathsf Q}$, $\mathsf K = ZA_{\mathsf K}$, $\mathsf V = ZA_{\mathsf V}$, output $\mathrm{softmax}(\mathsf
Q\mathsf K^\top/\sqrt h)\mathsf V$: each position averages all positions' values with weights from query--key similarity.
Cost $O(T^2h)$ in time and $O(T^2)$ in memory for the weights.
\iqlookfor{the formula, the scaling, and the quadratic cost.}
\end{solution}

\begin{solution}{iq:ml:sequence-models-on-order-books:3}
What the majority class share is; how the label is built (smoothing, horizon, dead band); how data were split (by day,
purged); results by day and instrument, with spreads; an IC or P\&L after costs; comparison with a linear model on
order-flow features.
\iqlookfor{class balance, label construction, split, and a baseline.}
\end{solution}

\begin{solution}{iq:ml:sequence-models-on-order-books:4}
The evidence (Kolm, Turiel and Westray; this chapter's learning curve) favours order-flow features at realistic data sizes:
they are stationary and already summarise the book. Raw levels can win with very large pooled datasets (Sirignano and
Cont). Start with the features; add raw inputs as extra channels.
\iqlookfor{data size as the deciding variable, with the published evidence.}
\end{solution}

\begin{solution}{iq:ml:sequence-models-on-order-books:5}
By whole sessions (days), in time order, with the test period after training and validation, purged at the edges; never
window by window at random; normalisation fitted on training days.
\iqlookfor{session-level splits and why windows leak.}
\end{solution}

\begin{solution}{iq:ml:sequence-models-on-order-books:6}
On short windows the most recent snapshots carry the signal; a TCN reads them directly and trains in parallel, while an
LSTM must carry them through its state and is slower to fit. An LSTM can do better when the relevant history is long and
irregular, with state to accumulate (a slowly building queue, a metaorder).
\iqlookfor{recency, training dynamics, and when recurrence pays.}
\end{solution}
